\chapter{Rectal Bleeding}

\section*{Definition}
Passage of blood from the rectum, which may be bright red (hematochezia, indicating distal GI source) or dark/maroon (suggesting proximal source).

\section*{When to See a Doctor}
See your doctor if:
\begin{itemize}
  \item Heavy rectal bleeding with dizziness or syncope, bleeding with abdominal pain, or any rectal bleeding in people over 50
\end{itemize}

\section*{How It Develops}
Bright red blood coating stool or on toilet paper typically originates from hemorrhoids (internal or external), anal fissures, or distal rectal pathology. Darker blood mixed with stool suggests a colonic source (diverticulosis, colitis, polyps, malignancy). Massive bright red bleeding with clots may indicate diverticular hemorrhage or angiodysplasia. Hemorrhoids are the most common cause but colorectal cancer must always be excluded.

\section*{Causes}
\begin{itemize}
  \item Internal or external hemorrhoids
  \item Anal fissure
  \item Diverticulosis (diverticular bleeding)
  \item Colorectal polyps or carcinoma
  \item Inflammatory bowel disease (ulcerative colitis, Crohn disease)
  \item Angiodysplasia (vascular malformation)
  \item Infectious colitis (E. coli, Shigella, Campylobacter)
  \item Radiation proctitis
  \item Ischemic colitis
  \item Rectal trauma or foreign body
\end{itemize}

\section*{Related Symptoms}
\begin{itemize}
  \item Blood in stool
  \item Anal pain
  \item Abdominal pain
  \item Constipation
  \item Weight loss
\end{itemize}


\section*{Medical Evaluation}
\begin{itemize}
  \item Digital rectal examination to palpate for masses, internal haemorrhoids, fissures, and prostate abnormalities; anoscopy for direct visualisation of the anal canal.
  \item Complete blood count to assess for anaemia; ferritin for iron deficiency; coagulation studies if bleeding history suggests a coagulopathy.
  \item Colonoscopy is selected when bleeding is unexplained or accompanied by alarm features such as anemia, weight loss, altered bowel habits, a relevant family history, or age-appropriate screening need; CT colonography may be considered when colonoscopy is incomplete or unsuitable.
  \item Referral to gastroenterology or colorectal surgery if colonoscopy is indicated, bleeding is recurrent or massive, or concerning polyps or masses are found.
\end{itemize}

\section*{Treatment Options}
\begin{itemize}
  \item For haemorrhoids: high-fibre diet, topical hydrocortisone 1\% cream or suppositories twice daily for seven days, warm sitz baths; rubber band ligation, sclerotherapy, or haemorrhoidectomy for persistent grade II--IV haemorrhoids.
  \item For anal fissure: stool softeners, topical nifedipine 0.3\% plus lidocaine 1.5\% ointment twice daily for six to eight weeks; lateral internal sphincterotomy for refractory cases.
  \item For diverticular bleeding: colonoscopic intervention (epinephrine injection, clipping) or interventional radiology with angiographic embolisation for active bleeding; surgical resection for uncontrolled haemorrhage.
  \item For inflammatory bowel disease: mesalamine derivatives, corticosteroids, or biologic therapy (anti-TNF agents) guided by gastroenterology.
\end{itemize}

\section*{Self-Care and Home Management}
\begin{itemize}
  \item Keep bowel movements soft: increase fibre to 25--35 g/day (psyllium husk, methylcellulose), drink eight to ten glasses of water daily, and use a stool softener if needed.
  \item Soak in a warm sitz bath for 10--15 minutes twice daily and after bowel movements to soothe the anal area and promote healing.
  \item Apply a topical haemorrhoidal cream (hydrocortisone 1\% or witch hazel) or use medicated wipes to reduce anal itching and inflammation.
  \item Avoid prolonged sitting on the toilet and straining during bowel movements, which worsen haemorrhoidal congestion.
\end{itemize}

\section*{References}
\begin{thebibliography}{99}
\bibitem{rectal_bleeding:ascrs-bleeding} Steele SR, et al. ``Clinical Practice Guideline for the Evaluation of Hematochezia.'' \emph{Dis Colon Rectum}. 2023;66(5):631-645.
\bibitem{rectal_bleeding:nejm-hematochezia} Wilkins T, et al. ``Evaluation and Management of Hematochezia.'' \emph{N Engl J Med}. 2024;390(12):1120-1131.
\bibitem{rectal_bleeding:harrison-gi-bleed} Longo DL, et al. \emph{Harrison\'s Principles of Internal Medicine}. 21st ed. McGraw-Hill; 2022. Chapter 324: Lower Gastrointestinal Bleeding.
\bibitem{rectal_bleeding:uptodate-lower-gi-bleed} Strate LL, et al. Approach to acute lower gastrointestinal bleeding in adults. In: UpToDate, Post TW (Ed), UpToDate, Waltham, MA; 2025.
\bibitem{rectal_bleeding:ascrs-hemorrhoids} Davis BR, et al. ``Clinical Practice Guideline for Hemorrhoidal Disease.'' \emph{Dis Colon Rectum}. 2023;66(6):760-773.
\bibitem{rectal_bleeding:cochrane-fibre-hemorrhoids} Alonso-Coello P, et al. ``Laxatives for the treatment of hemorrhoids.'' \emph{Cochrane Database Syst Rev}. 2022;(8):CD004649.
\end{thebibliography}
